% Generated draft: requires mathematical and source review.
\documentclass[11pt]{article}
\usepackage[margin=25mm]{geometry}
\usepackage{fontspec}
\setmainfont{Latin Modern Roman}
\usepackage{amsmath,amssymb,mathtools,unicode-math}
\setmathfont{Latin Modern Math}
\usepackage{enumitem,needspace,xurl,hyperref}
\setlength{\parindent}{0pt}
\setlength{\parskip}{5pt}
\setlength{\emergencystretch}{3em}
\newcommand{\N}{\mathbb{N}}
\newcommand{\Z}{\mathbb{Z}}
\newcommand{\Q}{\mathbb{Q}}
\newcommand{\R}{\mathbb{R}}
\newcommand{\C}{\mathbb{C}}
\newcommand{\F}{\mathbb{F}}
\newcommand{\A}{\mathbb{A}}
\newcommand{\PP}{\mathbb{P}}
\newcommand{\E}{\mathbb{E}}
\newcommand{\eps}{\varepsilon}
\newcommand{\dd}{\,\mathrm{d}}
\begin{document}
\section*{UnsolvedMath ID: 800680}
\textbf{Source problem code:} catalogue entry
\subsection{Definitions and mathematical statement}
Work on the complex projective plane \(\PP^2(\C)\). For a singular holomorphic foliation \(\mathscr F\) by curves, a leaf means a connected leaf in \(\PP^2(\C)\setminus\operatorname{Sing}(\mathscr F)\), equipped with its intrinsic Riemann surface structure. The geometric degree is the number of tangencies with a generic projective line, counted with multiplicity. An invariant line is tangent to the foliation at its regular points. We use the usual convention that a hyperbolic singularity has a local generating holomorphic vector field whose linear part has two nonzero eigenvalues with nonreal ratio.

A leaf \(L\) admits a Green function if, for every \(p\in L\), there exists a positive harmonic function \(G_p\colon L\setminus\{p\}\to(0,\infty)\) such that, in a local coordinate \(z\) centered at \(p\),
\[
G_p(z)=-\log|z|+O(1)\qquad(z\to0).
\]
This specifies existence of a positive Green function; no choice of normalization at infinity is required here.

\begin{enumerate}
\item[(i)] Let \(d\geq2\) be an integer, and let \(\mathscr F\) be a singular holomorphic foliation of geometric degree \(d\) on \(\PP^2(\C)\). Assume that the line at infinity is invariant and that every singularity of \(\mathscr F\) is hyperbolic. Does every leaf outside the source's exception, described as ``the compact curves,'' admit a Green function?

The compact-curve exception requires clarification. A compact invariant curve can contain singularities, whereas a leaf as defined above avoids singularities. Thus excluding compact leaves and excluding leaves contained in compact invariant curves are distinct formulations; the supplied record does not explicitly identify the intended one.

\item[(ii)] More generally, let \(d\geq2\) and assume \(\mathscr F\in\mathscr B_d\). Under the additional hypothesis that there exists a nonzero positive closed current directed by \(\mathscr F\), describe the function theory of the leaves in terms of Green functions.

In the projective-surface setting, a directed positive closed current means a positive current \(T\) of bidegree \((1,1)\), with \(T\neq0\) and \(\mathrm dT=0\), satisfying \(T\wedge\omega=0\) on regular coordinate neighborhoods where a holomorphic \(1\)-form \(\omega\) defines the foliation. The supplied record does not define \(\mathscr B_d\), so the admissible class and its additional hypotheses cannot be specified from this record. This part is a request for a description of leaf function theory, rather than an explicitly stated universal Green-function assertion.
\end{enumerate}
\subsection{Short English statement}
Dinh and Sibony ask whether leaves outside a compact-curve exception have Green functions for degree \(d\geq2\) foliations of \(\PP^2(\C)\) with invariant line at infinity and only hyperbolic singularities. They also ask for the Green-function theory of leaves when \(\mathscr F\in\mathscr B_d\) admits a nonzero positive closed directed current. The meaning of \(\mathscr B_d\) and the exact compact-curve exception require clarification.
\subsection{Sources}
\begin{itemize}
\item Local UnsolvedMath registry, permanent ID \(800680\), source code \texttt{catalogue entry}, title ``Holomorphic Foliation Theory --- Problem 7.'' The supplied record attributes the problem to Dinh and Sibony, with proposed year 2018, and labels it open. This is a supplied status label, not an independently verified mathematical assessment.
\item Dinh and Sibony, \emph{Some Open Problems on Holomorphic Foliation Theory} (2018), Problem 7. Supplied source URL: \url{https://arxiv.org/abs/1807.05640}. The supplied provenance records access on 2026-10-16 and extraction by \texttt{arxiv\_latex}; these provenance details are supplied metadata and were not independently verified for this draft.
\end{itemize}
\end{document}
